% 《大学物理学学习指导》第5版第10章“稳恒磁场”解题示例10.1--10.8
% 扫描PDF第150--156页，书内印刷142--148页。
% 原题参数和结构来自扫描逐页检查；推导独立重写，所有条目待独立二校。

\sourceexample{10.1}{142--143}{150--151}
\problemtext 一根均匀带电直线段$AB$长$b$，线电荷密度$\lambda$，
在垂直纸面的轴$O$周围以角速度$\omega$旋转，
$A$距$O$为$a$，求轴上磁感应强度$B_O$和磁矩$p_m$；
另求$a\gg b$时的近似式。
\workings
在$r$处取一电荷微段$dq=\lambda dr$，
旋转等效环电流为$dI=\omega dq/(2\pi)$。
一半径$r$的电流环在轴心的场为$dB=\mu_0dI/(2r)$，故
\[
\boxed{B_O=\frac{\mu_0\lambda\omega}{4\pi}
\ln\frac{a+b}{a}}.
\]
微元环的磁矩$d p_m=(\pi r^2)dI$，所以
\[
\boxed{p_m=\frac{\lambda\omega}{2}\int_a^{a+b}r^2dr
=\frac{\lambda\omega}{6}[(a+b)^3-a^3]}.
\]
当$b/a\ll1$，
$\ln(1+b/a)\simeq b/a$，立方差近似$3a^2b$，于是
\[
\boxed{B_O\simeq\frac{\mu_0\lambda\omega b}{4\pi a},
\qquad p_m\simeq\frac{\lambda\omega ba^2}{2}}.
\]
\crosscheck 对固定总电荷$q=\lambda b$，极限表达式正是半径$a$
的点电荷环绕旋转的中心场和磁矩。方向随$\lambda\omega$符号改变。

\sourceexample{10.2}{144}{152}
\problemtext 无限长实心圆柱导体半径$R$、电流$I$在截面均匀分布。
取沿轴向长$\ell=1$ m且在径向从轴线延伸至$2R$的矩形平面，
求穿过该矩形面的磁通量。
\workings
用安培定理得到
\[
B(r)=
\begin{cases}
\mu_0Ir/(2\pi R^2),&0<r<R,\\
\mu_0I/(2\pi r),&r>R.
\end{cases}
\]
沿矩形面取同向法线，场位于其法向，
\[
\Phi=\ell\left[\int_0^RB(r)dr+\int_R^{2R}B(r)dr\right]
=\frac{\mu_0I\ell}{2\pi}
\left(\frac12+\ln2\right).
\]
所以
\[
\boxed{\Phi=\frac{\mu_0I\ell}{4\pi}(1+2\ln2)}.
\]
\crosscheck 原书矩形从轴线延至$2R$，
不可遗漏圆柱内半径$0<r<R$的通量贡献。

\sourceexample{10.3}{144--145}{152--153}
\problemtext 同轴电缆由实心圆柱内导体半径$a$和空心外导体
内、外半径$b,c$组成，电流$I$在内导体轴向流出、
同量电流在外环截面均匀返回。求各半径的磁场大小。
\workings
取同轴圆形安培环路，
$B(2\pi r)=\mu_0I_{\rm enc}(r)$。
各区域包围电流分别为
\[
I_{\rm enc}=
\begin{cases}
I\,r^2/a^2,&0\le r<a,\\
I,&a\le r<b,\\
I[1-(r^2-b^2)/(c^2-b^2)],&b\le r<c,\\
0,&r\ge c.
\end{cases}
\]
因此
\[
\boxed{B(r)=\frac{\mu_0 I}{2\pi r}
\begin{cases}
r^2/a^2,&r<a,\\
1,&a\le r<b,\\
(c^2-r^2)/(c^2-b^2),&b\le r<c,\\
0,&r\ge c.
\end{cases}}
\]
\crosscheck 在$r=a,b,c$处相邻表达式连续；
电缆外部因往返电流相消而没有理想外磁场。

\sourceexample{10.4}{145--146}{153--154}
\problemtext 半径$R$、总电流$I$均匀分布的长圆柱形导体内
有半径$r$的偏心圆柱形空腔，腔轴$O'$离导体轴$O$距离$a>r$，
求$O,O'$两轴上的磁感应强度。
\workings
孔外材料的均匀电流密度为
$j=I/[\pi(R^2-r^2)]$。
用\emph{半径$R$、电流密度$j$的完整圆柱}
叠加\emph{半径$r$、电流密度$-j$的虚圆柱}模拟偏心空腔。
完整均匀载流圆柱在自己轴上的场为零，
虚圆柱在相距$a$处给出
\[
\boxed{B_O=\frac{\mu_0 j r^2}{2a}
=\frac{\mu_0 I r^2}{2\pi a(R^2-r^2)}}.
\]
在空腔轴$O'$处，虚圆柱自身场为零，
完整圆柱在距原轴$a$处给出
\[
\boxed{B_{O'}=\frac{\mu_0ja}{2}
=\frac{\mu_0Ia}{2\pi(R^2-r^2)}}.
\]
\crosscheck 两处磁场方位均垂直于轴心连线，
但用叠加表达式和电流方向判定矢量正负；
不能从电流总量$I$直接套用满截面$j=I/(\pi R^2)$。

\sourceexample{10.5}{146--147}{154--155}
\problemtext 无限长直导线$ab$通电$I_1$，在其一侧放置有限长直段$cd$通电$I_2$。
近端到长导线的垂距为$\ell_1$，短导线长度为$\ell_2$。
分别求$cd$沿垂直方向放置，以及与无限长直导线成角$\alpha$时的安培合力。
\workings
离长导线垂距为$x$处的磁感应强度
$B(x)=\mu_0I_1/(2\pi x)$。
第一种情形，导线元$d\ell=dx$，故
\[
\boxed{F_\perp=\frac{\mu_0 I_1I_2}{2\pi}
\ln\frac{\ell_1+\ell_2}{\ell_1}}.
\]
若短导线与长直导线夹角$\alpha$，
其径向投影$dx=d\ell\sin\alpha$，
积分界限$x\in[\ell_1,\ell_1+\ell_2\sin\alpha]$，
故
\[
\boxed{F_\alpha=
\frac{\mu_0I_1I_2}{2\pi\sin\alpha}
\ln\frac{\ell_1+\ell_2\sin\alpha}{\ell_1}}.
\]
\crosscheck 对$\alpha\to90^\circ$恢复垂直放置的结果；
对$\alpha\to0$作极限而非把$\sin\alpha=0$代入分母，
得到$F\to\mu_0I_1I_2\ell_2/(2\pi\ell_1)$。

\sourceexample{10.6}{147}{155}
\problemtext 显像管中电子在水平面内由南向北运动，
动能$E_k=1.2\times10^4\,\mathrm{eV}$；
地磁场的竖直向下分量$B_\perp=5.5\times10^{-5}$ T。
取$e=1.6\times10^{-19}$ C、$m_e=9.1\times10^{-31}$ kg。
求偏转方向、向心加速度，并求向北行进20 cm后横向偏移。
\workings
电子带负电，若速度向北、磁场向下，
$\bm v\times\bm B$对正电荷指西，
电子实际磁力则指\emph{东}，所以向东偏转。
动能换为SI单位$E_k=1.92\times10^{-15}$ J，
\[
v=\sqrt{\frac{2E_k}{m_e}}\approx6.50\times10^7\,\mathrm{m/s},
\quad
\boxed{a_n=\frac{evB_\perp}{m_e}\approx6.28\times10^{14}\,\mathrm{m/s^2}}.
\]
回旋半径
\[
R=\frac{m_ev}{eB_\perp}\approx6.72\,\mathrm m.
\]
北向位移$\ell=0.20$ m对应圆弧弦投影，
横向偏移
\[
\boxed{\delta=R-\sqrt{R^2-\ell^2}
\approx2.98\times10^{-3}\,\mathrm m}.
\]
\crosscheck $E_k=12$ keV时$v/c\approx0.22$，
本题依据原书采用非相对论近似；高精度需相对论修正。

\sourceexample{10.7}{148}{156}
\problemtext 测圆环形磁芯材料的相对磁导率：
匝数$N=200$、平均磁路长度$\ell=0.10$ m、
截面积$S=5.0\times10^{-5}\,\mathrm{m^2}$，
线圈电流$I=0.1$ A时截面磁通$\Phi=6.0\times10^{-5}$ Wb。
\workings
理想均匀环形磁路内，
\[
H=\frac{NI}{\ell}=200\,\mathrm{A/m},\qquad
B=\frac{\Phi}{S}=1.20\,\mathrm T.
\]
所以
\[
\boxed{\mu_r=\frac{B}{\mu_0H}
=\frac{\Phi\ell}{\mu_0NIS}
\approx4.77\times10^3}.
\]
\crosscheck 这是给定电流工作点的有效磁导率；
磁芯发生磁滞或饱和时$\mu_r$并非常数。

\sourceexample{10.8}{148}{156}
\problemtext 半径$R$的介质球具有空间均匀的磁化强度矢量$\bm M$，
无自由电流。求球心的磁感应强度$\bm B$及磁场强度$\bm H$。
\workings
均匀磁化球的体束缚电流密度
$\bm j_b=\nabla\times\bm M=\bm0$，
球面束缚面电流密度
$\bm K_b=\bm M\times\bm n$。
将球面划分为绕$\bm M$轴的细电流环，
用每个环对球心的轴向场积分，
可得
\[
\boxed{\bm B_O=\frac23\mu_0\bm M}.
\]
SI制中$\bm B=\mu_0(\bm H+\bm M)$，
从而
\[
\boxed{\bm H_O=-\frac13\bm M}.
\]
\crosscheck 结果与磁化球的退磁因子$N_d=1/3$相符；
$\bm B$沿磁化方向，而$\bm H$反向。
球心结果不依赖$R$，这是均匀磁化球的几何性质。
